% =====================================================================
%  CARNET D'ENTRAÎNEMENT — FICHE 5 (Bilans de matière)
%  THÈME 4 — Réactif limitant et tableau d'avancement
%  Fiche TUTO
% =====================================================================
\documentclass[11pt,a4paper]{article}
\usepackage[utf8]{inputenc}
\usepackage[T1]{fontenc}
\usepackage{lmodern}
\usepackage[french]{babel}
\usepackage{amsmath,amssymb,mathtools}
\usepackage{icomma}
\usepackage[version=4]{mhchem}
\usepackage{siunitx}
\sisetup{output-decimal-marker={,},per-mode=symbol}
\usepackage{xcolor}
\usepackage{colortbl}
\usepackage{tikz}
\usepackage{pgfplots}
\pgfplotsset{compat=1.18}
\usepackage[most]{tcolorbox}
\usepackage{enumitem}
\usepackage{array}
\usepackage{booktabs}
\usepackage{fancyhdr}
\usepackage[a4paper,top=1.5cm,bottom=1.4cm,left=1.8cm,right=1.8cm,headheight=24pt,headsep=6pt]{geometry}
\usetikzlibrary{arrows.meta,calc}

\definecolor{primary}{RGB}{142,93,168}
\definecolor{primarydark}{RGB}{82,45,108}
\definecolor{defcol}{RGB}{0,114,178}
\definecolor{thmcol}{RGB}{0,140,120}
\definecolor{attncol}{RGB}{200,40,55}
\definecolor{reacA}{RGB}{0,90,200}
\definecolor{reacB}{RGB}{210,30,40}
\definecolor{prodC}{RGB}{0,140,90}
\definecolor{prodD}{RGB}{198,93,9}
\newcommand{\nq}[1]{n_{\text{#1}}}

\pagestyle{fancy}\fancyhf{}
\fancyhead[L]{\small\textcolor{primarydark}{\textbf{Fiche 5} — Bilans de matière}}
\fancyhead[R]{\small\textcolor{primarydark}{\textsc{Tuto — Thème 4}}}
\fancyfoot[C]{\small\thepage}
\renewcommand{\headrulewidth}{0.8pt}
\renewcommand{\headrule}{\hbox to\headwidth{\color{primary}\leaders\hrule height \headrulewidth\hfill}}

\newtcolorbox{tutobox}[2][]{enhanced,breakable,boxrule=0.9pt,arc=2.5pt,
  left=3mm,right=3mm,top=2mm,bottom=2mm,colback=#2!6,colframe=#2,
  fonttitle=\bfseries,coltitle=white,title={#1}}

\setlength{\parindent}{0pt}
\setlength{\parskip}{2.3pt}

\begin{document}

\begin{center}
{\LARGE\bfseries\color{primarydark}Thème 4 — Réactif limitant \& avancement}\\[1pt]
{\color{primary}\rule{0.5\linewidth}{1.2pt}}\\[2pt]
{\itshape Qui s'épuise le premier ? C'est lui qui commande tout le bilan.}
\end{center}
\vspace{-3pt}

\begin{tutobox}[Réactif limitant vs réactif en excès]{defcol}
En pratique, on ne mélange pas les réactifs dans les proportions exactes. Le \textbf{réactif
limitant} est celui qui est \textbf{entièrement consommé le premier} : sa disparition
\emph{arrête} la réaction et fixe les quantités maximales de produits. Les autres réactifs, seulement
partiellement consommés, sont \textbf{en excès} : il en \emph{reste} à la fin.
\end{tutobox}

\begin{tutobox}[LE critère (à retenir absolument)]{thmcol}
Le limitant n'est \textbf{pas} le réactif en plus petite quantité ! On compare les rapports
$\dfrac{n_{\text{initial}}}{\text{coefficient}}$ ; le \textbf{plus petit rapport} désigne le limitant :
\[
\boxed{\;\text{limitant} = \text{réactif de plus petit } \dfrac{n_{i,0}}{\nu_i}\;}
\qquad\text{et}\qquad
\boxed{\;\xi_{\max}=\min_i\!\left(\dfrac{n_{i,0}}{\nu_i}\right)\;}
\]
\end{tutobox}

\begin{tutobox}[Le tableau d'avancement]{primary}
L'\textbf{avancement} $\xi$ (en mol) mesure « combien de fois » la réaction s'est produite. Pour
\ce{$a$A + $b$B -> $c$C + $d$D} :
\begin{center}\renewcommand{\arraystretch}{1.25}\small
\begin{tabular}{l|cccc}
 & $a\,\ce{A}$ & $b\,\ce{B}$ & $c\,\ce{C}$ & $d\,\ce{D}$\\ \hline
état initial ($\xi=0$) & $\nq{A,0}$ & $\nq{B,0}$ & $0$ & $0$\\
état $\xi$ & $\nq{A,0}-a\xi$ & $\nq{B,0}-b\xi$ & $c\xi$ & $d\xi$\\
état final ($\xi_{\max}$) & $\nq{A,0}-a\xi_{\max}$ & $\nq{B,0}-b\xi_{\max}$ & $c\,\xi_{\max}$ & $d\,\xi_{\max}$\\
\end{tabular}
\end{center}
Réactifs : on \emph{soustrait} $\nu_i\xi$ (ils décroissent). Produits : on \emph{ajoute} $\nu_i\xi$
(ils croissent). $\xi_{\max}$ est atteint quand \textbf{le premier réactif s'annule}.
\end{tutobox}

\begin{minipage}[t]{0.51\linewidth}
\vspace{0pt}
\begin{tutobox}[Lire le graphe $n(\xi)$]{attncol}
Pour $\ce{2A + 3B -> C + 6D}$ avec $\nq{A,0}=0{,}1$ et $\nq{B,0}=0{,}2$ :
les réactifs \emph{décroissent} (pentes $-2$ et $-3$), les produits \emph{croissent}
(pentes $+1$ et $+6$). \ce{A} s'annule le premier en $\xi_{\max}=0{,}05$ : \textbf{\ce{A} est limitant},
\ce{B} est en excès ($0{,}2-3\times0{,}05=0{,}05$~mol restantes).
\end{tutobox}
\end{minipage}\hfill
\begin{minipage}[t]{0.46\linewidth}
\vspace{0pt}
\centering
\begin{tikzpicture}
\begin{axis}[
   width=6.9cm,height=5.2cm,
   xlabel={\scriptsize avancement $\xi$ (mol)},
   ylabel={\scriptsize $n$ (mol)},
   xmin=0,xmax=0.06,ymin=0,ymax=0.32,
   axis lines=left,
   xtick={0,0.025,0.05},ytick={0,0.1,0.2,0.3},
   tick label style={font=\tiny},label style={font=\tiny},
   legend style={font=\tiny,at={(1.02,1)},anchor=north west,draw=black!25},
   clip=false,
]
  \addplot[reacA,line width=1pt] coordinates {(0,0.10)(0.05,0)}; \addlegendentry{$\nq{A}$}
  \addplot[reacB,line width=1pt] coordinates {(0,0.20)(0.05,0.05)}; \addlegendentry{$\nq{B}$}
  \addplot[prodC,line width=1pt] coordinates {(0,0)(0.05,0.05)}; \addlegendentry{$\nq{C}$}
  \addplot[prodD,line width=1pt] coordinates {(0,0)(0.05,0.30)}; \addlegendentry{$\nq{D}$}
  \draw[densely dashed,black!55] (axis cs:0.05,0)--(axis cs:0.05,0.30);
  \node[black!70,font=\tiny,anchor=south] at (axis cs:0.05,0.305){$\xi_{\max}$};
  \fill[reacA] (axis cs:0.05,0) circle (1.6pt);
\end{axis}
\end{tikzpicture}
\end{minipage}

\begin{tutobox}[Méthode complète du bilan avec limitant]{thmcol}
\textbf{1.} Convertir les données en moles. \;
\textbf{2.} Calculer $n_{i,0}/\nu_i$ pour \emph{chaque} réactif $\to$ le plus petit $=$ limitant $=\xi_{\max}$. \;
\textbf{3.} Produits formés : $n=\nu_{\text{produit}}\times\xi_{\max}$. \;
\textbf{4.} Excès restant : $n_{\text{reste}}=n_{i,0}-\nu_i\,\xi_{\max}$.
\end{tutobox}

\end{document}
